\documentclass[crop=false]{standalone}
% --- ПРЕАМБУЛА ДЛЯ ПОДДЕРЖКИ РУССКОГО ЯЗЫКА И МАТЕМАТИКИ ---
\usepackage[T2A]{fontenc}
\usepackage[utf8]{inputenc}
\usepackage[russian]{babel}
\usepackage{amsmath}
\usepackage{geometry}
\usepackage{amssymb}
\usepackage{graphicx}
\usepackage{float}
\usepackage{import}
\geometry{a4paper, top=2cm, bottom=2cm, left=2cm, right=2cm}
\newcommand{\sidecomment}[1]{\marginpar{\raggedright\scriptsize\itshape #1}}
\newcommand{\iu}{\mathrm{i}}

\newcommand{\Def}{%
    \text{Def}%
    \hspace{-0.85em}% Сдвиг назад ровно на середину (подберите под шрифт)
    \raisebox{0.1ex}{/}% Черта
    \hspace{1em}% Возврат к концу слова + небольшой отступ
}

\renewcommand{\Re}{\operatorname{Re}}
\renewcommand{\Im}{\operatorname{Im}}

\ifstandalone
    \newcommand{\lecturebreak}{} % В отдельном файле лекции разрыв не нужен
\else
    \newcommand{\lecturebreak}{\clearpage} % В полной книге - начинаем с новой страницы
\fi

\newcounter{lecturenumber} % Создаем счетчик для нумерации лекций
\newcommand{\lecture}[2]{
    \lecturebreak
    \refstepcounter{lecturenumber} % Увеличиваем счетчик на 1
    \par\vspace{2em}\noindent % Отступ сверху и убираем абзацный отступ
    {\Large\bfseries % Делаем заголовок крупным и жирным
    Лекция \thelecturenumber: #1 % Выводим "Лекция N: Тема"
    \hfill % Растягиваем пробел, чтобы дата была справа
    \large\textit{#2}} % Выводим дату курсивом
    \par\vspace{1em} % Небольшой отступ снизу
    \addcontentsline{toc}{section}{Лекция \thelecturenumber: #1} % Добавляем в оглавление
    \par
}

\newcommand{\TwoCols}[2]{%
  \par\noindent % Начать с новой строки без отступа
  \begin{minipage}[t]{0.48\textwidth}
    #1
  \end{minipage}% <-- Важный '%' для удаления пробела
  \hfill % Растягивающийся пробел между колонками
  \begin{minipage}[t]{0.48\textwidth}
    #2
  \end{minipage}% <-- Важный '%' для удаления пробела
  \par%
}

\pagestyle{plain} % Включаем нумерацию страниц\
\begin{document}
% fixme везде съехало \hat{\dot q_i}
\lecture{}{29.05}
\section{Преобразования Лежандра}
\[X(x_1, \dots, x_n) \qquad H_X \not \equiv 0\]
где $H_X$ — гессиан $X$
Новые переменные $y_1, \dots, y_n$:
\[y_i = y_i(x_1, \dots, x_n) = \frac{\partial X}{\partial x_i}\]
Вопрос, можно ли построить $Y(y_1, \dots, Y_n)$ , что старые переменные будут 
\[x_i = \frac{\partial Y}{\partial y_i}\]
Можно, это называется преобразованием Лежандра

\noindent Замечание: функция $X$ по отношению к такой замене координат является производящей функцией

\noindent \keypoint
% fixme ставить \hat x сомнительно считаю, имеет смысл переделать, пока запишу так
\[X(x_1, x_2) = \ln x_1 + x_2^2 \qquad x_i > 0\]
\[y_1 = \frac{1}{x_1} \qquad y_2 = 2 x_2\]
\[x_1 = \frac{1}{y_1} = \frac{\partial Y}{\partial y_1} \qquad x_2 = \frac{1}{2} y_2 = \frac{\partial Y}{\partial y_2}\]
\[Y = \ln y_1 + \frac{y_1^2}{4}\]
\[\hat Y(x_1, x_2) = Y(y_1(x_1, x_2), y_2(x_1, x_2)) = \ln \frac{1}{x_1} + x_2^2\]
\[\hat X = \ln \frac{1}{y_1} + \frac{y_2^2}{4} \qquad \hat x_i = \hat x_i(y_1, \dots ,y_n)\]
\[Y = \sum_{i=1}^{n} \hat x_i y_i - \hat X\]
\[\frac{\partial Y}{\partial y_i} = \frac{\partial}{\partial y_i} \left( \sum_{j=1}^{n} \hat x_i y_i - \hat X\right) = \hat x_i + \sum_{j=1}^{n} y_i \frac{\partial \hat x_i}{\partial y_i} - \sum_{j=1}^{n} \frac{\partial X}{\partial x_j} \frac{\partial \hat x_j}{\partial y_i} = \hat x_i\]
\[H_Y \not \equiv 0 \quad \boxed{Y = \sum_{i=1}^{n} \hat x_i y_i - \hat X} \iff \boxed{X = \sum_{i=}^{n} \hat y_i x_i - \hat Y}, \quad H_X \not \equiv 0\]
\begin{proof}
    симметрично, очевидно, можно подставить вместо 
    \[\hat Y = \widehat{\left(\sum_{i=1}^{n} \hat x_i y_i - \hat X\right)}\] 
\end{proof}
\section{Преобразование Гамильтона}

\keypoint \[L = T - \Pi = \frac{mv^2}{2} = \frac{m \ddot x^2}{2} + \frac{m \ddot y^2}{2} + \frac{m \ddot z^2}{2} - \Pi(x,y,z)\]
\[m \bar v = \begin{pmatrix}
    m v_x \\ m v_y \\ m v_z
\end{pmatrix} = \begin{pmatrix}
    m \dot x  \\ m \dot y \\ m \dot z
\end{pmatrix}\qquad \frac{\partial L}{\partial \dot x_i} = m \dot x\]
\Def Обобщённый импульс 
\[p_i  = \frac{\partial L}{\partial \dot q_i}\]
Замечание: конкретно в этом примере получилось, что обобщённые импульсы имеют тот же физический смысл, что и обычные
\Def Ф-я Гамильтона (Гамильтониан)$H$ — это преобразование Лежандра функции лагранжа по обобщённым скоростям
\[H(p_1, \dots, p_n, q_1, \dots q_n, t)\]
\[p_i = \frac{\partial L}{\partial \dot q_i} \qquad L(\dot q_1, \dots, \dot q_n, q_1, \dots, q_n , t)\]
\[H(p_1, \dots, p_n, \dot q_1, \dots, \dot q_n ,t) = \sum_{i=1}^{n} p_i \hat{\dot q_i} - \hat L \qquad H = \sum_{i=1}^{n} p_i \hat{\dot q_i} - \hat L \qquad L = \sum_{i=1}^{n} \hat p_i \dot q_i - \hat H \]
\[\dot q_i = \frac{\partial H}{\partial p_i}\]
\section{Уравнения Гамильтона}
\[d H = \sum_{i=1}^{n} \frac{\partial H}{\partial p_i} d p_i + \sum_{i=1}^{n} \frac{\partial H}{\partial q_i} d q_i + \frac{\partial H}{\partial t} d t\]
\[d H = d \left( \sum_{i=1}^{n} p_i \hat {\dot q_i} - \hat L\right) = \sum_{i=1}^{n} \hat{\dot q_i} d p_i + \sum_{i=1}^{n} p_i d \hat{\dot q_i} - \sum_{i=1}^{n} \underbracket{\frac{\partial L}{\partial \dot q_i}}_{p_i}  d \dot q_i - \sum_{i=1}^{n} \frac{\partial L}{\partial q_i} d q_i - \frac{\partial L}{\partial t} d t\]
% fixme почему d \hat{\dot q_i} = d \dot q_i
% fixme посмотреть у Дениса всё ли верно в конспекте
\[H = \sum_{i=1}^{n} p_i \hat{\dot q_i} - \hat L \qquad L = \sum_{i=1}^{n} \hat p_i \dot q_i - \hat H\]
\[\dot q_i = \frac{\partial H}{\partial p_i} \qquad \frac{\partial H}{\partial t} = - \frac{\partial L}{\partial t} \qquad \frac{\partial H}{\partial q_i} = - \frac{\partial L}{\partial q_i}\]
Правая часть уравнения Лагранжа нуль
\[\frac{d }{d t} \underbracket{\frac{\partial L}{\partial \dot q_i}}_{p_i} - \frac{\partial L}{\partial q_i} = 0\]
Уравнения Гамильтона
\[\boxed{\dot q_i = \frac{\partial H}{\partial p_i} \qquad \dot p_i = -\frac{\partial H}{\partial q_i}}\]
\Def Гамильтонова система — механическая система, описываемая уравнениями Гамильтона
\keypoint Консервативная системы всегда Гамильтоновы
\section{Физический смысл Гамильтониана}
\[L = T - \Pi \qquad L = L_2 + L_1 + L_0\]
\[\hat H = \sum_{i=1}^{n} \hat p_i \dot q_i - L = \sum_{i=1}^{n} \frac{\partial L}{\partial \dot q_i} \dot q_i - L = \sum_{i=1}^{n} \underbrace{\frac{\partial L_2}{\partial \dot q_i} \dot q_i}_{2 L_2} + \sum_{i=1}^{n} \underbrace{\frac{\partial L_1}{\partial \dot q_i} \dot q_i}_{L_1} + \sum_{i=1}^{n} \underbrace{\frac{\partial L_0}{\partial \dot q_i} \dot q_i}_0 - L_2 - L_1 - L_0 = L_2 - L_0\]
Равенства в силу теоремы Эйлеры об однородной функции
\[H = L_2 - L_0 = T_2 - T_0 + \Pi\]
В Склерономных системах $T_0 = 0, T_1 = 0$
\[H = T_2 + \Pi = T + \Pi\]
То есть в Консервативных системах гамильтониан равен полной энергии мех. системы
\end{document}